\chapter{The Jarvis-Operas Package}
\label{ch:jarvis-operas}

\newthought{A registry of reusable operators.} \Operas\ is the core plugin package
behind two features you have already met: named operators for \yamlkey{Operas.Modules}
(Chapter~\ref{ch:operas}) and qualified functions callable inside every expression
(Chapter~\ref{ch:expressions}). This chapter explains the registry itself: what lives
in it, how Workers discover it, and how to put your own functions in.

\section{What the registry holds}
\label{sec:operas-registry}

Functions live under namespaces --- \code{math.add}, \code{helper.eggbox2d},
\code{user.my\_loglike} --- populated from three sources:

\begin{itemize}
  \item \textbf{built-ins and catalog functions} shipped with the package;
  \item \textbf{persisted user operators} --- functions you register once with
        \Operas's own tools, stored so any process can rediscover them;
  \item \textbf{entry points} --- any installed Python package advertising
        \code{jarvis\_operas.core} or \code{jarvis\_operas.user} entry points.
\end{itemize}

Inspect an installation from the command line:

\begin{terminal}
\begin{jcmd}
Jarvis operas list            # every registered name, by namespace
Jarvis operas info math.add   # signature and documentation
\end{jcmd}
\end{terminal}

\section{The two consumption paths}
\label{sec:operas-consumption}

\paragraph{As a workflow step.} A registered name is a valid
\yamlkey{Operas.Modules[].operator} --- tried after dotted-import resolution
(Section~\ref{sec:opera-resolution}):

\begin{yamlwin}
Operas:
  Modules:
    - name: EggBox
      operator: helper.eggbox2d
      output: [{name: z, entry: z}]
\end{yamlwin}

\paragraph{Inside expressions.} Every registered function is callable by its qualified
name in any expression field --- likelihood terms, opera inputs, \yamlkey{Dump}
variables, \yamlkey{selection}, \sampler{AdaptiveBridson} targets:

\begin{yamlwin}
Sampling:
  LogLikelihood:
    - name: LogL
      expression: user.external_loglike(x, y)
\end{yamlwin}

There is deliberately \emph{no} function-declaration block in the task card: the card
names functions; the registry provides them.

\section{Discovery and the spawn boundary}
\label{sec:operas-discovery}

Each Worker is a fresh \code{spawn} process, so each discovers the registry
independently at startup --- persisted operators and entry points are found, a local
parsing/numerical snapshot is built, and compiled expressions call the snapshotted
functions directly with no per-Sample registry lookups. Discovery is demand-gated:
cards using only built-in functions never initialize the package at all.

\begin{jwarn}
The spawn boundary has one consequence worth engraving: \textbf{a function registered
ad hoc in your launching Python process does not exist in the Workers.} Anything the
card references must be rediscoverable --- persist it into the registry or install it
as an entry point. ``Works in my notebook, \code{NameError} in the Worker'' is this
rule, every time.
\end{jwarn}

\section{Registering your own}
\label{sec:operas-registering}

Two supported routes:

\begin{description}[style=nextline, leftmargin=1.2em]
  \item[Persist with \Operas's tools] --- register the function once via the package's
        own \textsc{cli}/\textsc{api} (see its documentation); it lands in the
        persisted user store under \code{user.<name>} and every Worker on the machine
        finds it. Quickest for personal, single-machine work.
  \item[Package an entry point] --- expose the function from an installed package
        under the \code{jarvis\_operas.user} (or \code{core}) entry-point group. The
        right answer for anything shared: versioned, \code{pip}-installable on the
        cluster, visible to \cli{Jarvis operas list} everywhere.
\end{description}

\begin{jtip}
Keep registered functions \emph{pure and picklable-free}: take numbers, return
numbers (or a mapping, for operator use); load heavy assets lazily inside the module,
not at import. That keeps Worker startup fast and makes the same function usable both
as a workflow operator and inside a likelihood expression.
\end{jtip}
